\chapter{Ретроспективный аудит обменного KO6-каркаса}

Перед построением нового Clifford-линейного цикла необходимо проверить,
не был ли его существенный каркас уже строго получен в ранних томах.
Настоящий гейт перечитывает доказанные участки версий III--V и строгую
ветвь RPFT с точки зрения обменной вещественной формы Тёплица.

\section{Что уже было доказано}

\subsection{Обязательное KO6-удвоение}

Гейт \path{version3_real_bimodule_square_gate.tex} доказал, что
недублированный квадрат имеет $J\gamma=\gamma J$, а KO6 требует
$J\gamma=-\gamma J$. Минимальное исправление имеет универсальный вид
\begin{equation}
 H=H_p\oplus H_p^c,
 \qquad
 \gamma=\gamma_p\oplus(-\gamma_p),
 \qquad
 J(\xi,\eta)=(\overline\eta,\overline\xi).
 \label{eq:v5-retro-universal-ko6-doubling}
\end{equation}
Тот же каркас затем был явно проверен на семейном колчане размерности 18,
относительной цепи Пати--Салама размерности 20 и аффинном модуле
размерности 300.

Это ровно алгебраическая часть текущей обменной Real-формы: две
ориентированные копии не являются независимыми объектами, а составляют
одну $J$-орбиту.

\subsection{Обмен ориентированного дифференциала}

В аффинном KO6-гейте уже получено точное равенство
\begin{equation}
 JdJ^{-1}=d^\dagger.
 \label{eq:v5-retro-j-d-exchange}
\end{equation}
По форме оно совпадает с обменом операторов
\begin{equation}
 T_+=S\otimes q_0,
 \qquad
 T_-=S^*\otimes\overline{q_0}.
\end{equation}
Следовательно, для следующего шага не требуется придумывать новое
конечномерное particle--antiparticle удвоение. Совместимость $J$,
градуировки и сопряжённых стрелок уже является повторно проверенной леммой
проекта.

\subsection{Словарь следов и мер}

Гейт \path{version3_orbit_measure_pfaffian_gate.tex} строго разделил
полный бозонный спектральный след и фермионный Pfaffian half-count.
Гейт \path{version4_pfaffian_eta_orientation_gate.tex} показал, что два
противоположных reduced Pfaffian-знака превращаются в положительный
квадрат модуля после полного KO6-спаривания.

Из этого следует важное ограничение нового чтения:
\begin{equation}
 \text{сокращение полного следа или Pfaffian-фазы}
 \ne
 \text{вычисление класса в }KO_6.
 \label{eq:v5-retro-trace-not-kclass}
\end{equation}
Поэтому нулевой суммарный комплексный индекс не опровергает обменный
вещественный класс, но и нормированный вес $1/7$ сам по себе его не
доказывает.

\subsection{Запрет извлекать ориентацию только из KO6}

Гейт \path{version5_oriented_height_hodge_ko6_gate.tex} построил две
неэквивалентные KO6-совместимые высоты и доказал: вещественность и обмен
$J$ не выбирают знак ориентации. В текущей ветви знак поэтому должен
приходить не из $J$, а из уже объявленных хопфовых данных
\begin{equation}
 c_1(L)=+1,
 \qquad
 c_1(L^*)=-1,
\end{equation}
и из выбора символов $z$ и $z^{-1}$ в расширении Тёплица.

\section{Что ранняя программа не даёт}

Все перечисленные KO6-модули конечномерны. Для любого квадратного
конечномерного оператора обычный индекс равен нулю. Они проверяют
вещественность, градуировку, условие первого порядка и следовые
нормировки, но не являются фредгольмовыми представителями граничного
класса Тёплица.

Особенно важно не переносить сюда старую попытку строгой RPFT вывести
топологический коэффициент через семейство голономий $\theta:0\to\pi$.
Файл \path{rigorous/22_spectral_flow_derivation.py} прямо фиксирует
отсутствие пересечения нуля, а
\path{rigorous/21_pi_coefficient_derivation.md} заключает, что индекс
этого семейства равен нулю. Его последующие аргументы используют
геометрическую и планковскую нормировку, а не доказанный индекс. Поэтому
эта архивная ветвь служит отрицательным контролем, а не основанием нового
класса.

\begin{theorem}[Ретроспективное разделение труда]
Ранние тома уже дают полный алгебраический KO6-каркас обменной пары:
сопряжённые копии, противоположные градуировки, антилинейный обмен $J$ и
правило $JdJ^{-1}=d^\dagger$. Они не дают аналитического класса, поскольку
не содержат бесконечномерного фредгольмова представителя, Clifford-сдвига
степени шесть и отображения сравнения с двумя комплексными ориентациями.
\end{theorem}

\section{Как меняется следующий шаг}

Первоначально планировалось сразу строить полный $Cl_{0,6}$-линейный
оператор. Ретроспектива показывает, что это слишком широкий гейт. Его
алгебраическая половина уже закрыта, а единственный неизвестный узел ---
отображение сравнения
\begin{equation}
 KO_6(\mathbb C_{\mathbb R})
 \longrightarrow
 K_0(\mathbb C\oplus\mathbb C)
 \label{eq:v5-retro-real-complex-comparison}
\end{equation}
для комплексификации обменной вещественной формы.

Следует сначала определить образ генератора. Успехом является
антидиагональная ориентированная пара
\begin{equation}
 1\longmapsto(-1,+1)
 \quad\text{с точностью до общего знака}.
 \label{eq:v5-retro-required-antidiagonal}
\end{equation}
После умножения на коэффициентный проектор ранга пятнадцать она даст
$(-15,+15)$. Если же каноническое отображение даёт диагональную пару или
обнуляет требуемую компоненту, кандидат текущей ветви закрывается до
любых физических интерпретаций.

Clifford/Bott-стабилизация должна при этом входить в аналитический
$KKO$-цикл. Её нельзя читать как восемь, шестнадцать или иное число новых
физических фермионов: ранние гейты уже доказали, что прямое добавление
вспомогательного KO6-модуля в физическое конечное пространство меняет
состав модели.

\begin{nogocriterion}[Запрет повторного строительства конечной геометрии]
Следующий гейт не должен снова перебирать конечномерные матрицы $J$ и
$\gamma$ или добавлять новый particle--antiparticle сектор. Он считается
содержательным только при явном вычислении отображения
\eqref{eq:v5-retro-real-complex-comparison} и его совместимости с
граничной картой Тёплица.
\end{nogocriterion}

\begin{protocol}
\begin{enumerate}
 \item ретроспективно проверенные тома: \textbf{III, IV и V};
 \item строгая архивная ветвь RPFT: \textbf{проверена};
 \item готовый обменный KO6-каркас: \textbf{найден};
 \item новое конечномерное удвоение: \textbf{не требуется};
 \item полный след или Pfaffian как тест $KO_6$-класса:
 \textbf{запрещён};
 \item старая RPFT-ветвь спектрального потока: \textbf{индекс ноль};
 \item готовый $Cl_{0,6}$-индексный представитель: \textbf{не найден};
 \item требуемый новый результат:
 \textbf{вещественно-комплексное отображение сравнения};
 \item критерий успеха: \textbf{образ генератора $(-1,+1)$};
 \item критерий остановки: \textbf{диагональный или нулевой образ};
 \item физическое замыкание: \textbf{не объявлено};
 \item следующий гейт:
 \textbf{Bott/комплексификационное отображение обменной пары}.
\end{enumerate}
\end{protocol}

Машинный сертификат сохранён в
\path{s2t_v5_real_toeplitz_cross_tome_reuse_audit_gate_results.json}.